\documentclass[conference]{IEEEtran}
\IEEEoverridecommandlockouts

\usepackage{cite}
\usepackage{amsmath,amssymb,amsfonts}
\usepackage{graphicx}
\usepackage{textcomp}
\usepackage{xcolor}
\usepackage{booktabs}
\usepackage{url}
\usepackage[T1]{fontenc}
\usepackage{balance}
\usepackage{hyperref}
\hypersetup{
    colorlinks=true,
    linkcolor=black,
    urlcolor=blue,
    citecolor=black,
}

\def\BibTeX{{\rm B\kern-.05em{\sc i\kern-.025em b}\kern-.08em
    T\kern-.1667em\lower.7ex\hbox{E}\kern-.125emX}}

\begin{document}

\title{How Much of Turkish Wordle Is Luck? \\
A Multi-Baseline Skill/Luck Split Using Letter Frequency
and Information Theory}

\author{\IEEEauthorblockN{Bar\i{}\c{s} Cem Bayburtlu}
\IEEEauthorblockA{
\textit{Department of Software Engineering} \\
\textit{\c{C}ankaya University} \\
Ankara, Turkey \\
c2228009@student.cankaya.edu.tr
\thanks{Source code, data, and figures are available at
\url{https://github.com/byigitt/what-percentage-wordle-is}.}
}
}

\maketitle

\begin{abstract}
This paper measures how much of winning at Turkish Wordle
(\textit{wordleturkce.bundle.app}) is due to luck and how much is due to
skill. We first extract the full word lists from the live game: 2{,}936
unique answer words and 5{,}500 valid guesses. We then simulate four
player models on every puzzle in the pool. A player that only follows
the colored feedback (green, yellow, gray), but uses no word-selection
strategy, already wins $96.97\%$ of the games. This shows that feedback
alone carries a large amount of information. However, this player is
allowed to know every word in the dictionary, which is not realistic for
most humans. We therefore add a vocabulary-limited human model, where
the player only knows the top $K$ common words, and an entropy-greedy
bot in the style of the \textit{3Blue1Brown} lesson. The entropy bot wins
every puzzle with $3.59$ guesses on average. With $K{=}1000$ (covering
around $48\%$ of daily puzzles), the measured split is about $34\%$ luck
and $66\%$ skill. Inside that skill share, vocabulary explains about
$95\%$ and strategy explains about $5\%$. We also report four other
baselines, where the headline luck share changes from $0.1\%$ to $97\%$.
The main message is that a skill/luck number for Wordle is only
meaningful after the baseline player is made explicit.
\end{abstract}

\begin{IEEEkeywords}
Wordle, information theory, Shannon entropy, Turkish language, word games,
skill versus luck, simulation, computational linguistics
\end{IEEEkeywords}

\section{Introduction}

Wordle was made by Josh Wardle in 2021 and was later bought by The New
York Times. In the game, the player has six tries to find a hidden
five-letter word. After each guess, the game returns one color for each
letter position: green means the right letter is in the right place,
yellow means the right letter is in the wrong place, and gray means the
letter is not in the word. Repeated letters need one extra counting rule,
because the same letter can appear more than once in a guess or answer.

The game became popular because the rules are simple, but the decisions
are not. A player must choose between common words, words with many
different letters, and words that separate the remaining candidates well.
This makes Wordle a useful small test case for information theory,
language structure, and human word knowledge. Many local versions also
appeared, including the Turkish-language
\textit{wordleturkce.bundle.app}, which uses its own list of Turkish
five-letter words.

A common question among players is: how much of a win is luck, and how
much is skill? The question sounds simple, but it is not well defined.
``Luck'' could mean that the daily word was easy. It could also mean
that the player happened to know the word, or that the feedback happened
to remove many wrong answers. In the same way, ``skill'' could mean
knowing many words, choosing a good opener, or using feedback correctly.
A percentage such as ``$X\%$ luck and $Y\%$ skill'' therefore needs a
baseline: a clear description of what the no-skill or lower-skill player
can do.

Most Wordle papers and solver discussions focus on finding better
guesses. That is important, but it answers a different question. Solver
papers compare algorithms. This paper asks how the answer changes when
we compare a strong solver to several weaker players, including a
human-like player whose vocabulary is limited. This is the main part of
the work that is different from previous Wordle studies: we do not only
ask which solver is best; we also measure how much of the gap can be
explained by vocabulary, feedback, and strategic word choice.

The Turkish version is also worth studying on its own. Turkish word
forms and letter positions are not the same as English. Five-letter
Turkish words often place vowels and consonants in different positions
than English words, and common Turkish words do not have the same letter
frequency pattern. For that reason, results from English Wordle should
not be copied directly to Turkish Wordle without checking the data.

This paper answers the question with five steps:
\begin{enumerate}
\item Extract the answer pool and guess dictionary from the shipped
JavaScript bundle (Section~\ref{sec:data});
\item Build four player models, from random guessing to an
information-theoretic bot (Sections~\ref{sec:models}, \ref{sec:entropy});
\item Add a human-like player whose vocabulary is limited to the top
$K$ common words (Section~\ref{sec:human});
\item Compute the skill/luck split under five baselines, so the
baseline sensitivity is visible (Section~\ref{sec:decomposition}); and
\item Study per-word difficulty and the word families that remain hard
even for the entropy bot (Section~\ref{sec:difficulty}).
\end{enumerate}

\section{Related Work}

Wordle solvers got a lot of attention soon after the game came out.
Sanderson's video~\cite{3b1b_wordle} made the entropy-greedy approach
widely known. It uses Shannon's information theory~\cite{shannon1948}:
each turn the player picks the guess with the highest expected
information gain, $H(g \mid C) = -\sum_p P(p \mid g, C) \log_2 P(p \mid g, C)$.
Later work showed that the truly optimal English bot needs about
$3.42$ guesses on average and wins every game by doing full game-tree
search~\cite{bertsimas2022wordle}. A recent peer-reviewed paper by
Aladaileh et al.~\cite{aladaileh2026wordle} confirms again that entropy
beats simple letter-distribution heuristics on English Wordle.

For Turkish Wordle, Akbulut and Senturk~\cite{akbulut2026turkish}
already made the first careful solver study. They compare LSTM, MaxEnt,
and rule-based methods on $1{,}000$ Turkish target words. They also show
that linguistic features, especially dense CVCVC templates and low word
frequency, help predict word difficulty. This paper uses that result as
background, but it asks a different research question. We do not try to
beat other solvers. Instead, we ask how much of a win remains after we
choose a baseline player. We also use the full $2{,}936$-word solution
pool from the deployed game and add a human-like player with limited
vocabulary.

The vocabulary-limited idea has a close relative in the HSW model of
Qiu and Zhong~\cite{qiu2024hsw}, who use a vocabulary-population (VP)
component based on a normal-distribution assumption to simulate
heterogeneous English players. Our model is different: we do not assume
any distribution. Instead we rank words by how often the deployed
bundle picks them as the daily answer.

For games in general, the skill-vs-luck question has been studied more
formally. The most rigorous work is on poker~\cite{poker_luck}, which
uses a predictive-likelihood argument. We follow a simpler
counterfactual idea from Mauboussin~\cite{mauboussin2012success}: a
game outcome counts as skill if the skilled agent wins where the
baseline agent would lose, and as luck if both win or both lose. For
English Wordle, Dilger~\cite{dilger2023wordle} applies information
theory to a large sample of player data and uses the share of
one-guess solves to estimate a cheating rate. Our work is different:
we report the skill/luck split under several human baselines, not a
single score against one bot.

This leaves a gap that the present paper targets. We are not aware of
earlier Wordle work that combines three pieces in one experiment: a
multi-baseline skill/luck split, a vocabulary-limited human baseline,
and the full deployed Turkish Wordle solution pool. This matters because
Turkish is not just English with different letters. It is agglutinative,
and many five-letter Turkish words follow CVCVC-like shapes. Therefore
the letter-position counts, good openers, and hard word families must be
checked on Turkish data directly.

\section{Data and Methodology}
\label{sec:data}

\subsection{Word Lists}

The game runs fully in the browser. The complete word list ships inside
one JavaScript file at
\url{https://wordleturkce.bundle.app/static/main.2023-04-23.js}. We
find two arrays in it:

\begin{itemize}
\item \texttt{La}: 11{,}470 entries, but only 2{,}936 unique. The daily
answer is $La[i \bmod 11470]$, where $i$ is the day index. The same
word can appear many times---for example \texttt{buruk} and
\texttt{selam} each appear 8 times, while rare words appear once. We
read this repetition as a commonality signal: the curator made
common words more likely to be the daily answer.
\item \texttt{Ta}: 5{,}531 entries, 5{,}499 unique. These count as
valid guesses but never appear as the answer.
\end{itemize}

After we drop duplicates, the full guess dictionary has 5{,}500 unique
words. The daily answer comes from the 2{,}936 unique solutions, with
the $La$ multiplicities used as weights.

\subsection{Feedback Function}

We implement the standard Wordle feedback rule, including the case
with repeated letters. It is a deterministic function
$\mathrm{feedback}: \Sigma^5 \times \Sigma^5 \to \{0,1,2\}^5$, where
0, 1, and 2 stand for gray, yellow, and green. We assign greens first.
Then we walk through the non-green positions of the guess from left to
right and mark a yellow when the answer still has an unmatched copy of
that letter. We encode each five-cell pattern as a base-3 integer in
$[0, 242]$ so that we can index it directly.

\subsection{Letter Frequency Analysis}

Over the weighted 11{,}470-entry solution distribution, the five most
common letters are A ($13.04\%$), E ($8.52\%$), \.I ($6.90\%$),
K ($6.81\%$), and R ($6.09\%$). Together they cover $41\%$ of all
letter occurrences. By position:

\begin{table}[h]
\centering
\caption{Most frequent letters by position (top 5 each).}
\label{tab:positional}
\begin{tabular}{cl}
\toprule
Position & Top 5 letters (with counts) \\
\midrule
1 & K (1336), S (1121), T (841), B (751), M (714) \\
2 & A (3199), E (1881), \.I (1142), O (938), U (782) \\
3 & R (1435), L (1183), K (799), N (757), M (747) \\
4 & A (2308), \.I (1587), E (1530), I (868), U (843) \\
5 & K (1367), N (1189), R (1070), A (1052), E (971) \\
\bottomrule
\end{tabular}
\end{table}

This pattern matches Turkish phonotactics. Turkish words use a
canonical $(C)V(C)$ syllable template~\cite{goksel2005turkish}, and
CVCVC is one of the most common five-letter shapes. Vowels are mostly
in positions 2 and 4 and consonants in positions 1, 3, and 5.

\section{Player Models}
\label{sec:models}

We use four player models. All of them follow the standard rules: up
to six guesses, and each guess must be a valid five-letter word from
the dictionary.

\subsection{Pure Random ($\mathcal{P}_{\mathrm{rand}}$)}
At each turn this player picks a random word from the 5{,}500-word
dictionary and ignores all feedback. This is the zero-information
baseline. The expected win rate in six guesses is
$1 - (1 - 1/5500)^6 \approx 0.109\%$.

\subsection{Consistent Random ($\mathcal{P}_{\mathrm{cons}}$)}
This player keeps the set of candidates that are still consistent with
the feedback so far. At each turn it picks a random word from that
set. It follows the rules perfectly but uses no strategy beyond that.

\subsection{Frequency Bot ($\mathcal{P}_{\mathrm{freq}}$)}
At each turn this bot recomputes positional letter frequencies on the
current candidate set. Each candidate $w$ receives a score
$\mathrm{score}(w) = \sum_i \mathbb{1}[w[i] \notin w[0{:}i]] \cdot
\mathrm{posfreq}[i, w[i]]$, which sums the contributions of the
\emph{distinct} letters only. The bot picks the candidate with the
highest score.

\subsection{Vocabulary-Limited Human ($\mathcal{P}_{\mathrm{hum}, K}$)}
\label{sec:human}
This model represents a casual human player by a vocabulary size $K$.
We rank words by how often \texttt{La} repeats them (the commonality
proxy), and break ties alphabetically. The top $K$ are the player's
vocabulary. At each turn, the player picks a random word from the
intersection of the remaining candidates and the vocabulary. If that
intersection is empty, the player picks a random word from the
vocabulary as a forced ``probe'' guess.

\section{Entropy-Greedy Bot}
\label{sec:entropy}

To get a tight skill ceiling we add an entropy-greedy bot
$\mathcal{P}_{\mathrm{ent}}$ in the style of \cite{3b1b_wordle}. For a
guess $g$ and candidate set $C$, the partition that $g$ creates on $C$
is $\{C_p : p \in \{0,\ldots,242\}\}$ where
$C_p = \{a \in C : \mathrm{feedback}(g,a) = p\}$. The expected
information gain is the entropy of this partition:
\begin{equation}
H(g \mid C) = -\sum_{p} \frac{|C_p|}{|C|} \log_2 \frac{|C_p|}{|C|}.
\end{equation}
The bot picks $g^\star = \arg\max_g H(g \mid C)$ and, when several
guesses are tied, prefers a guess that is itself in $C$.

\subsection{Implementation}

A direct implementation needs $O(|G| \cdot |C|)$ feedback evaluations
per turn, which is too slow when we simulate thousands of puzzles. To
speed this up we precompute a pattern matrix
$M \in \{0,\ldots,242\}^{|G| \times |A|}$ once. We store the encoded
feedback for every guess--answer pair as a \texttt{uint8}. The matrix
takes about $8\,$MB and is built in about $0.5\,$s using vectorized
NumPy~\cite{harris2020numpy} operations that work in parallel over
letters.

To check the fast feedback function we compared it with a plain Python
reference on $5{,}000$ random word pairs. We found 0 mismatches.

\subsection{Results}

The best entropy opener is \texttt{merak} with $H = 5.81$ bits. The top
entropy openers include \texttt{merak}, \texttt{kenar}, \texttt{kamer},
\texttt{kelam}, \texttt{kalem}, \texttt{malik}, \texttt{karne},
\texttt{katre}, \texttt{keman}, and \texttt{karni}. This list shows that mid-frequency
liquid consonants (R, L, N, M) can be more useful than the very frequent
letter A, which dominated the frequency-score top picks.

On all 2{,}936 puzzles, the entropy bot wins
\textbf{$100.00\%$} of the time, with $3.59$ guesses on average. It
never reaches turn six. The turn histogram is
$\{1: 1, 2: 84, 3: 1221, 4: 1455, 5: 175\}$.

\section{Simulation Results}
\label{sec:results}

\begin{table}[t]
\centering
\caption{Performance of all four player models on the 2{,}936-word
solution pool (uniform sampling).}
\label{tab:players}
\begin{tabular}{lrrr}
\toprule
Player & Win rate & Avg.\ guesses & Opener \\
\midrule
$\mathcal{P}_{\mathrm{rand}}$ (pure random) & 0.10\% & 4.33 & --- \\
$\mathcal{P}_{\mathrm{hum},1000}$ (K=1000) & 33.92\% & 3.84 & --- \\
$\mathcal{P}_{\mathrm{cons}}$ (consistent) & 96.97\% & 4.21 & --- \\
$\mathcal{P}_{\mathrm{freq}}$ (frequency) & 99.18\% & 3.77 & \texttt{salik} \\
$\mathcal{P}_{\mathrm{ent}}$ (entropy) & \textbf{100.00\%} & \textbf{3.59} & \texttt{merak} \\
\bottomrule
\end{tabular}
\end{table}

Table~\ref{tab:players} shows the main numbers. Fig.~\ref{fig:turns}
shows the turn distribution for each model. The largest gap is between
$\mathcal{P}_{\mathrm{rand}}$ ($0.1\%$) and
$\mathcal{P}_{\mathrm{cons}}$ ($96.97\%$). Once a player uses the
feedback correctly, most of the win rate is already recovered. Strategy
on top of that adds a smaller gain.

\begin{figure}[t]
\centering
\includegraphics[width=\columnwidth]{figures/turn_breakdown.png}
\caption{Turn counts for the four player models on 2{,}936 puzzles.
Pure random almost always runs out of guesses. Consistent random and
the frequency bot have long tails. The entropy bot is tightly grouped
and never goes past five guesses.}
\label{fig:turns}
\end{figure}

\subsection{Vocabulary Sensitivity}

\begin{table}[t]
\centering
\caption{Win rates of the vocabulary-limited human model and the
share of daily puzzles its vocabulary covers, as a function of $K$.}
\label{tab:human}
\begin{tabular}{rrrr}
\toprule
$K$ & Uniform win \% & Weighted win \% & Coverage \% \\
\midrule
200  & 6.81  & 10.92 & 10.92 \\
500  & 17.00 & 25.74 & 25.79 \\
1000 & 33.92 & 47.39 & 47.59 \\
1500 & 50.72 & 67.00 & 67.50 \\
2000 & 67.13 & 83.68 & 84.93 \\
2500 & 83.58 & 94.46 & 96.20 \\
2936 & 97.21 & 97.20 & 100.00 \\
\bottomrule
\end{tabular}
\end{table}

Table~\ref{tab:human} and Fig.~\ref{fig:k} show how the win rate
depends on the vocabulary size $K$. The empirical fit
\begin{equation}
P(\mathrm{win} \mid K) \approx \mathrm{coverage}(K) \cdot 0.97
\label{eq:vocab}
\end{equation}
holds for every $K$ we tested, within sampling noise. Equation
\ref{eq:vocab} gives a direct interpretation: \emph{knowing the word is
almost necessary and sufficient to win}. If the day's word is in the
player's vocabulary, consistent-random play wins with probability about
$0.97$. Otherwise the player usually loses.

\begin{figure}[t]
\centering
\includegraphics[width=\columnwidth]{figures/k_vs_winrate.png}
\caption{Win rate of the vocabulary-limited human model as a function
of vocabulary size $K$. The weighted curve is almost identical to the
vocabulary coverage curve, which matches Eq.~\ref{eq:vocab}. The
shaded band is the plausible range for typical casual players.}
\label{fig:k}
\end{figure}

\section{Skill/Luck Decomposition}
\label{sec:decomposition}

We use a counterfactual decomposition. The entropy bot's wins count as
\emph{luck} when the baseline player would have won too, and as
\emph{skill} when the baseline player would have lost. For a baseline
$\mathcal{P}_b$ with win rate $w_b$ and a ceiling
$\mathcal{P}_{\mathrm{ent}}$ with win rate $w_c = 1$, the luck share is
$w_b / w_c$ and the skill share is $(w_c - w_b) / w_c$.

\begin{table}[t]
\centering
\caption{Skill/luck split under five different baselines. The result
changes by almost two orders of magnitude.}
\label{tab:decomp}
\begin{tabular}{lrrr}
\toprule
Baseline $\mathcal{P}_b$ & $w_b$ & Luck \% & Skill \% \\
\midrule
Pure random              & 0.10\%  & 0.1\%  & 99.9\% \\
Human ($K{=}500$)        & 17.00\% & 17.0\% & 83.0\% \\
\textbf{Human ($K{=}1000$)} & \textbf{33.92\%} & \textbf{33.9\%} & \textbf{66.1\%} \\
Human ($K{=}2000$)       & 67.13\% & 67.1\% & 32.9\% \\
Consistent random        & 96.97\% & 97.0\% & 3.0\%  \\
\bottomrule
\end{tabular}
\end{table}

Table~\ref{tab:decomp} and Fig.~\ref{fig:decomp} show the split. The
headline answer depends entirely on the baseline we pick. We argue
that the most realistic baseline for an average human is
$\mathcal{P}_{\mathrm{hum},K \approx 1000}$, since informal estimates
put active Turkish five-letter vocabulary around this size for casual
players. Under this baseline,
\textbf{$\approx{}66\%$ of Turkish Wordle wins come from skill and
$\approx{}34\%$ come from luck.}

\begin{figure}[t]
\centering
\includegraphics[width=\columnwidth]{figures/skill_decomposition.png}
\caption{Skill and luck shares of entropy-bot wins under five
baselines. The answer to ``how much of Wordle is luck'' moves from
$0.1\%$ to $97\%$ just based on which baseline we pick.}
\label{fig:decomp}
\end{figure}

\subsection{Skill is Mostly Vocabulary}

Equation~\ref{eq:vocab} also tells us that the skill share itself
splits into two parts. Going from $\mathcal{P}_{\mathrm{hum},1000}$ to
$\mathcal{P}_{\mathrm{cons}}$ (perfect vocabulary, no strategy) adds 63
points of win rate ($33.92\%$ to $96.97\%$). Going from
$\mathcal{P}_{\mathrm{cons}}$ to $\mathcal{P}_{\mathrm{ent}}$ (strategy
on top of perfect vocabulary) adds only 3 points more ($96.97\%$ to
$100.00\%$). So out of the 66-point skill gap, vocabulary explains
about $95\%$ and strategic word selection explains about $5\%$.

This differs from the usual story in popular Wordle posts, which focuses
on opener choice and information-theoretic optimization. Our result
suggests that \emph{lexical knowledge} matters much more than strategic
sophistication for who wins more often.

\section{Per-Word Difficulty}
\label{sec:difficulty}

Fig.~\ref{fig:diff} shows how many turns the entropy bot needs to
solve each of the 2{,}936 puzzles. The most common turn count is 4.
The bot never reaches turn 6. The puzzles that take 5 turns (175 words)
are the structurally hardest ones.

\begin{figure}[t]
\centering
\includegraphics[width=\columnwidth]{figures/difficulty_hist.png}
\caption{Turn counts of the entropy-greedy bot on all 2{,}936
solutions. No puzzle needs 6 turns; the hardest 175 take 5.}
\label{fig:diff}
\end{figure}

When we look inside the 5-turn group we see a common pattern: large
families of look-alike words that differ in only one position, for
example \texttt{KAPAK / KAÇAK / KAVAK / KAZAK}. Once the candidate set
shrinks to such a family, even an optimal bot has to test the
remaining words one by one because each guess tells it about one
position only. This is an irreducible source of variance. No
strategy can remove it, and an ideal human would face the same
problem.

This kind of luck is \emph{structural}, and is different from the
player-level luck we measured earlier. Even a perfect-information
player has to spend more turns when the day's word lives in a dense
look-alike family.

\section{Discussion}
\label{sec:discussion}

\subsection{Definitional Sensitivity}

The main methodological point of this work is that the popular
``$X\%$ luck, $Y\%$ skill'' framing is not well defined without a
baseline. The same data, looked at under five reasonable baselines,
gives skill shares anywhere between $3\%$ and $99.9\%$. We therefore
think single-number claims of this kind should be avoided in both
popular and academic discussion of game design.

\subsection{Vocabulary as the Dominant Skill}

The near-identity between vocabulary coverage and win rate
(Eq.~\ref{eq:vocab}) points to an under-discussed finding: Wordle is
mostly a vocabulary test. Strategy helps, but its marginal effect is
small. A player who wants a higher win rate may gain more by learning
more words than by changing to a more complex opening word. This
matches independent observational work by Liang et
al.~\cite{liang2024biases}, who show that real Wordle players narrow
candidates using semantic, orthographic, and morphological
similarities, not information-theoretic optimality.

\subsection{Information-Theoretic Bound}

The Shannon entropy of the weighted daily distribution is 11.36 bits.
The theoretical ceiling per guess is $\log_2 243 = 7.92$ bits. The
entropy bot pulls out about $3.16$ bits per guess on average (the
same as needing $3.59$ guesses to remove $11.36$ bits of
uncertainty). That is around $40\%$ of the per-guess ceiling. The gap
has two sources: (i) the greedy choice is not the same as full
lookahead, and (ii) every guess must be a real word, not just any
five-character string.

\subsection{Limitations}

The vocabulary model uses bundle multiplicity as a proxy for
commonality. A more rigorous human model would calibrate against an
external Turkish corpus and would also include imperfect feedback
tracking, especially for words with repeated letters. Stronger
solver models, such as the Bertsimas--Paskov-style full-tree
search~\cite{bertsimas2022wordle}, would raise the skill ceiling
slightly but would not change the overall split.

\section{Conclusion}

We have presented a multi-baseline analysis of skill versus luck in
Turkish Wordle, built on a full extraction of the deployed game's
solution and dictionary lists. Four player models were simulated on
all 2{,}936 puzzles: the entropy-greedy ceiling wins every game, while
a realistic vocabulary-limited human (K=1000) wins about $34\%$.
Under this baseline, the skill share of wins is about $66\%$, and
about $95\%$ of that is vocabulary while only about $5\%$ is strategy.
We stress that the headline split is very sensitive to the baseline
and we report a five-baseline table to back up this point.

All code, simulation outputs, intermediate data, and reproduction
steps are available as open source at the project repository.

\section*{Reproducibility}

The source code, extracted word lists, full simulation outputs, and
reproduction steps are available in the project's open-source
repository. The whole pipeline runs end-to-end in under three minutes on
a normal laptop, using only NumPy~\cite{harris2020numpy} and
Matplotlib~\cite{hunter2007matplotlib}.

\balance
\begin{thebibliography}{00}

\bibitem{shannon1948}
C.~E.~Shannon, ``A mathematical theory of communication,''
\textit{Bell System Technical Journal}, vol.~27, no.~3, pp.~379--423, 1948.

\bibitem{3b1b_wordle}
G.~Sanderson (3Blue1Brown), ``Solving Wordle using information theory,''
YouTube video, 2022.\\
\url{https://www.youtube.com/watch?v=v68zYyaEmEA}

\bibitem{bertsimas2022wordle}
D.~Bertsimas and A.~Paskov,
``An Exact Solution to Wordle,''
\textit{Operations Research}, vol.~73, no.~3, pp.~1384--1394, 2025.

\bibitem{mauboussin2012success}
M.~J.~Mauboussin, \textit{The Success Equation: Untangling Skill and
Luck in Business, Sports, and Investing}. Harvard Business Review Press,
2012.

\bibitem{poker_luck}
S.~D.~Levitt and T.~J.~Miles, ``The Role of Skill Versus Luck in Poker:
Evidence from the World Series of Poker,''
\textit{Journal of Sports Economics}, vol.~15, no.~1, pp.~31--44, 2014.

\bibitem{wordle_nyt}
J.~Wardle, ``Wordle,'' originally published 2021;
\url{https://www.nytimes.com/games/wordle/index.html}.

\bibitem{wordleturkce}
``Wordle Türkçe,'' \url{https://wordleturkce.bundle.app/}, accessed 2026.

\bibitem{aladaileh2026wordle}
T.~Aladaileh, D.~Stephens, M.~Alqaisi, and C.~Wu,
``Solving Wordle Using Information Theory,''
\textit{Northeast Journal of Complex Systems}, vol.~8, no.~1, art.~6, 2026.\\
\url{https://orb.binghamton.edu/nejcs/vol8/iss1/6/}

\bibitem{harris2020numpy}
C.~R.~Harris \textit{et al.}, ``Array programming with NumPy,''
\textit{Nature}, vol.~585, no.~7825, pp.~357--362, 2020.

\bibitem{hunter2007matplotlib}
J.~D.~Hunter, ``Matplotlib: A 2D graphics environment,''
\textit{Computing in Science \& Engineering}, vol.~9, no.~3, pp.~90--95,
2007.

\bibitem{liang2024biases}
G.~Liang, A.~Kabbara, C.~Liu, R.~Luo, K.~Kim, and M.~Guerzhoy,
``Semantic, Orthographic, and Phonological Biases in Humans' Wordle
Gameplay,'' arXiv preprint arXiv:2411.18634, 2024.\\
\url{https://arxiv.org/abs/2411.18634}

\bibitem{goksel2005turkish}
A.~Göksel and C.~Kerslake,
\textit{Turkish: A Comprehensive Grammar}. London: Routledge, 2005.

\bibitem{akbulut2026turkish}
S.~Akbulut and \.{I}.~F.~Senturk, ``An analytical study of artificial
intelligence methods in the Turkish Wordle game: LSTM, entropy, and
rule-based approaches,'' \textit{Entertainment Computing}, vol.~57,
art.~101117, May 2026, doi: 10.1016/j.entcom.2026.101117.

\bibitem{qiu2024hsw}
T.~Qiu and Y.~Zhong, ``Simulation and analysis of Wordle based on the HSW
model,'' in \textit{Frontiers in Artificial Intelligence and Applications},
vol.~376, IOS Press, 2024, pp.~167--174.

\bibitem{dilger2023wordle}
J.~P.~Dilger, ``Wordle: A microcosm of life. Luck, skill, cheating,
loyalty, and influence!'' arXiv preprint arXiv:2309.02110, 2023.\\
\url{https://arxiv.org/abs/2309.02110}

\end{thebibliography}

\end{document}
